\section{尝试策略}

\subsection{策略一：上下文裁剪}

对历史对话进行裁剪，只保留与当前轮请求直接相关的内容。该策略可以降低 prefill token 数，从而减少 TTFT，但需要避免删掉影响任务理解的关键信息。

\subsection{策略二：稳定 Prompt 前缀}

将系统提示词、工具说明和固定格式模板放在稳定前缀中，减少每轮请求中前缀内容的变化。稳定前缀有利于 prefix KV cache 命中，适合多轮对话 Agent 的连续请求场景。

\subsection{策略三：会话级 KV Cache 复用}

为同一个会话维护可复用的 KV cache，使新一轮请求只对新增用户输入和必要增量上下文进行 prefill。该策略对降低多轮 TTFT 最直接，但需要管理 cache 生命周期、显存占用和会话淘汰策略。

\subsection{策略四：工具调用后置或异步化}

对于不影响首轮响应的工具调用，尝试后置或异步执行，先让模型返回初始响应或状态提示。该策略可以降低端到端 TTFT，但只适用于工具结果不是首 token 必需条件的场景。

\subsection{策略五：请求调度优化}

通过限制最大 batch 等待时间、区分短上下文和长上下文请求、为交互式请求设置更高优先级等方式降低排队时间。该策略主要缓解高并发下 TTFT 抖动。
